\documentclass[11pt]{amsart}
\usepackage{amsmath,amssymb,amsthm}
\usepackage[margin=1.12in]{geometry}
\usepackage[hidelinks]{hyperref}

\newtheorem{theorem}{Theorem}
\newtheorem{lemma}{Lemma}
\newtheorem{corollary}{Corollary}

\newcommand{\R}{\mathbb R}
\newcommand{\B}{B_2^n}
\newcommand{\voln}[1]{|#1|_n}
\newcommand{\volnm}[1]{|#1|_{nm}}

\title{Equality in the Polar Schneider Collection Inequality}
\author{arxiv\_automatic\_math run fa\_banach\_001}
\date{2026-07-03}

\begin{document}
\maketitle

\section*{Source and Status}

This packet concerns Julian Haddad, Dylan Langharst, Galyna V. Livshyts, and
Eli Putterman, \emph{On the polar of Schneider's difference body},
arXiv:2503.06191.

Status: \texttt{full\_solution\_likely\_valid} for the equality-condition gap
in the source paper's Theorem 11.

\section*{Source Question}

For a collection \(\mathcal K=(K_0,\ldots,K_m)\) of convex bodies in \(\R^n\),
the source defines
\[
  D^m(\mathcal K)
  =\{(x_1,\ldots,x_m):K_0\cap(K_1+x_1)\cap\cdots\cap(K_m+x_m)\neq\varnothing\}.
\]
Theorem 11 of the source paper states that, if all \(K_i\) have center of mass
at the origin and \(\voln{K_i^\circ}=\voln{K_0^\circ}\), then
\[
  \volnm{D^{m,\circ}(\mathcal K)}
  \prod_{i=1}^m\voln{K_i}
  \leq
  \voln{\B}^m\,\volnm{D^{m,\circ}(\B)}.
\]
It further states that equality is necessary only if each \(K_i\) is a centered
ellipsoid, and sufficient if all \(K_i\) are the same centered ellipsoid.  The
authors explicitly note a gap: they do not know whether equality can occur when
all \(K_i\) are centered ellipsoids but at least one differs from the others.

\section*{Idea of the Proof}

The only missing case is a tuple of centered ellipsoids.  Write
\[
  h_{E_i}(u)^2=u^TQ_i u
\]
with \(Q_i\) positive definite.  Equal polar volumes are exactly the condition
that \(\det Q_i\) is independent of \(i\).  The polar volume of
\(D^m(E_0,\ldots,E_m)\) is an integral of
\(\exp(-h_{D^m(E_0,\ldots,E_m)})\).  The elementary Laplace representation of
\(e^{-\sqrt{s}}\) writes each factor as a positive mixture of Gaussians.

For each fixed set of mixture parameters \(t_0,\ldots,t_m>0\), the inner
Gaussian integral is controlled by the determinant of a block matrix.  The
source paper already computes this determinant in its functional section:
\[
  \det M
  =
  \prod_{i=0}^m\det(t_iQ_i)\,
  \det\left(\sum_{i=0}^m(t_iQ_i)^{-1}\right).
\]
Minkowski's determinant inequality gives a pointwise lower bound for this
determinant, with strict equality unless all \(Q_i\) are equal.  Since the
mixture is positive, the strict pointwise inequality integrates to strict
volume loss for every non-identical ellipsoid tuple.

\section*{Main Theorem}

\begin{theorem}
Let \(m,n\in\mathbb N\).  Under the hypotheses of Theorem 11 of Haddad,
Langharst, Livshyts, and Putterman, equality holds if and only if
\[
  K_0=K_1=\cdots=K_m
\]
is a centered ellipsoid.
\end{theorem}

\begin{proof}
The source paper already proves the inequality, proves equality for identical
centered ellipsoids, and proves that equality can occur only when every
\(K_i\) is a centered ellipsoid.  It remains to rule out equality for different
centered ellipsoids.

Let \(E_i=A_i\B\) be centered ellipsoids and put \(Q_i=A_iA_i^T\), so
\[
  h_{E_i}(u)=\sqrt{u^TQ_i u}.
\]
The condition \(\voln{E_i^\circ}=\voln{E_0^\circ}\) is equivalent to
\[
  \det Q_i=\det Q_0=:\delta
  \qquad(0\leq i\leq m).
\]
Indeed \(\voln{E_i^\circ}=\voln{\B}/\sqrt{\det Q_i}\).

For \(z=(z_1,\ldots,z_m)\in(\R^n)^m\), the support-function formula in the
source paper gives
\[
  h_{D^m(E_0,\ldots,E_m)}(z)
  =
  \sqrt{\left(\sum_{j=1}^mz_j\right)^TQ_0\left(\sum_{j=1}^mz_j\right)}
  +
  \sum_{i=1}^m\sqrt{z_i^TQ_i z_i}.
\]
Using the standard polar-volume formula in dimension \(N=nm\),
\[
  \volnm{D^{m,\circ}(E_0,\ldots,E_m)}
  =
  \frac{1}{(nm)!}
  \int_{(\R^n)^m}e^{-h_{D^m(E_0,\ldots,E_m)}(z)}\,dz.
\]

We use the positive Gaussian-mixture identity
\[
  e^{-r}
  =
  \frac{1}{\sqrt{2\pi}}
  \int_0^\infty
    t^{-3/2}\exp\left[-\frac12\left(tr^2+t^{-1}\right)\right]\,dt
  \qquad(r\geq0).
\]
Applying this identity to the \(m+1\) square-root terms in the support function
and using Fubini, the volume integral is a positive average, over
\(t_0,\ldots,t_m>0\), of Gaussian integrals of the form
\[
  \int_{(\R^n)^m}
  \exp\left[
    -\frac12\left(
      t_0\left(\sum_{j=1}^mz_j\right)^TQ_0\left(\sum_{j=1}^mz_j\right)
      +\sum_{i=1}^m t_i z_i^TQ_i z_i
    \right)
  \right]\,dz.
\]
The corresponding block matrix \(M(t,Q)\) has diagonal blocks
\(t_iQ_i+t_0Q_0\) and off-diagonal blocks \(t_0Q_0\).  By the determinant
calculation in Proposition 30 of the source paper, applied to the matrices
\(t_iQ_i\),
\[
  \det M(t,Q)
  =
  \prod_{i=0}^m\det(t_iQ_i)\,
  \det\left(\sum_{i=0}^m(t_iQ_i)^{-1}\right).
\]
The inner Gaussian integral is \((2\pi)^{nm/2}\det M(t,Q)^{-1/2}\).

By Minkowski's determinant inequality for positive-definite matrices,
\[
 \det\left(\sum_{i=0}^m t_i^{-1}Q_i^{-1}\right)^{1/n}
 \geq
 \sum_{i=0}^m t_i^{-1}\det(Q_i^{-1})^{1/n}
 =
 \delta^{-1/n}\sum_{i=0}^m t_i^{-1}.
\]
Equivalently,
\[
  \det\left(\sum_{i=0}^m(t_iQ_i)^{-1}\right)
  \geq
  \delta^{-1}\left(\sum_{i=0}^m t_i^{-1}\right)^n.
\]
Therefore
\[
  \det M(t,Q)
  \geq
  \delta^m\left(\prod_{i=0}^m t_i^n\right)
  \left(\sum_{i=0}^m t_i^{-1}\right)^n.
\]
This lower bound is exactly the value of \(\det M(t,Q)\) when all
\(Q_i\) are replaced by a single positive-definite matrix \(Q\) with
\(\det Q=\delta\).

The equality condition in Minkowski's determinant inequality says that equality
holds only when the matrices \(Q_i^{-1}\) are pairwise proportional.  Since all
\(\det Q_i\) are equal, this means \(Q_0=\cdots=Q_m\).  Hence, if the ellipsoid
tuple is not identical, then for every \(t_i>0\),
\[
  \det M(t,Q)
  >
  \delta^m\left(\prod_{i=0}^m t_i^n\right)
  \left(\sum_{i=0}^m t_i^{-1}\right)^n,
\]
and the corresponding Gaussian integral is strictly smaller than for the
identical ellipsoid tuple.  Since the mixing measure above is positive, the
strict inequality survives integration.  Thus
\[
  \volnm{D^{m,\circ}(E_0,\ldots,E_m)}
  <
  \volnm{D^{m,\circ}(E,\ldots,E)}
\]
whenever the \(E_i\) are not all equal and \(E\) is any centered ellipsoid with
\(\det(E)=\det(E_i)\).

Finally,
\[
  \prod_{i=1}^m\voln{E_i}=\voln{E}^m,
\]
so a non-identical ellipsoid tuple gives a strictly smaller left-hand side than
the identical ellipsoid tuple.  The identical tuple is an equality case by the
source paper.  Therefore equality in Theorem 11 occurs exactly for identical
centered ellipsoids.
\end{proof}

\begin{corollary}
The earlier \(m=1\) partial packet is a special case: if \(E_0,E_1\) are
centered ellipsoids with \(\voln{E_0^\circ}=\voln{E_1^\circ}\), equality in the
\(m=1\) collection inequality holds if and only if \(E_0=E_1\).
\end{corollary}

\section*{Limitations}

This solves the equality gap in the source paper's polar Schneider collection
inequality.  It does not solve Schneider's original higher-order conjecture
that \(\volnm{D^m(K)}\voln{K}^{-m}\) is minimized by ellipsoids for
\(n\geq3\) and \(m\geq2\).

\section*{Novelty Check}

Local run indexes were searched for \texttt{2503.06191}, ``Polar Schneider'',
``Theorem 11'', ``collection equality'', ``same centered ellipsoid'', and
``Rogers-Brascamp-Lieb-Luttinger equality''.  The only local packet was the
earlier \(m=1\) partial result, now superseded by this packet.

On 2026-07-03, bounded web searches for the exact title and phrases including
``Polar Schneider inequality equality collection ellipsoids'', ``same centered
ellipsoid'', and ``Rogers-Brascamp-Lieb-Luttinger Schneider equality'' surfaced
the source arXiv page and no separate answer to the equality gap.  Novelty
confidence is moderate: the proof uses elementary tools and a determinant
calculation already present in the source paper's functional section, so the
observation may be accessible to the authors or experts.

\section*{Human Review Recommendation}

Review as a likely valid full solution to the source equality-condition gap.
The key checks are:

\begin{itemize}
\item the positive Gaussian-mixture identity for \(e^{-r}\);
\item the determinant formula for the block Gaussian matrix after replacing
  the source matrices by \(t_iQ_i\);
\item the strict equality condition in Minkowski's determinant inequality.
\end{itemize}

\begin{thebibliography}{9}

\bibitem{HLLP}
J. Haddad, D. Langharst, G. V. Livshyts, and E. Putterman,
\emph{On the polar of Schneider's difference body},
arXiv:2503.06191, 2025.

\bibitem{AGA}
S. Artstein-Avidan, A. Giannopoulos, and V. D. Milman,
\emph{Asymptotic Geometric Analysis, Part I},
Mathematical Surveys and Monographs 202,
American Mathematical Society, 2015.

\end{thebibliography}

\end{document}
